% Page budget: 1.55 pages | Owns: negation, OBC derivation, additional-state scope, implementation, and true-parameter condition.
% WRITING GUIDE
% Purpose: Explain negation, derive the final state-level orthogonal-by-construction map, and delimit its interpretability guarantee.
% Sources: README "Source map per subsection", rows IV-A, IV-B and IV-C; mandatory papers and the reference gate as in the README.
% Research-plan anchor: Preserve Aspect 3's goal of protecting physical interpretability; state clearly how the final construction differs from the proposed regularisation.
% Current decisions: Reduced identifiable coordinates, data-derived reference tuples, a tangent-only one-step physical-state projection, per-epoch refresh, and no affine comparison arm.
% Choices to justify: Protected parameter coordinates, reference-set construction, projected rows, rank tolerance, exclusion of the affine offset, refresh frequency, and treatment of additional states.
% Cautions: docs/orthogonal-projection-plan.md describes an older soft-penalty route; one-step state orthogonality is not trajectory-output orthogonality or guaranteed true recovery.
% Planned figures: New geometric projection figure with a physical and additional-state partition inset; show the remaining uniqueness limitation as clearly as the protected direction.
% Section is finished when: The protected space, coordinate frame, gradient treatment, added-state scope, and true-recovery condition are explicit.
%
% MUST ESTABLISH (PROPOSED, cycle 11, inferred from README ownership row IV, source-map rows IV-A to IV-D,
%   THESIS-RESULTS step 5; the file held no block)
%  IV-A 1. The network of Section III can cancel a change of theta_base to first order (negation, Kessels
%          Sec. 5.3.1.3; Gyorok 2026 Ex. 2, 2025 Ex. 1); parameter directions = sensitivity in the ten
%          identifiable training coordinates xi; on the gantry a negating write is an inertial, damping or
%          stiffness force whose direction depends on Y.
%  IV-B 2. Adapted OBC (Gyorok 2026 Eqs. 8, 10, Lemma 3, Eq. 13): regressor = sensitivity at an expansion
%          point (Gyorok 2025 Sec. 4), data-derived tuples at xbar = 0, six physical rows, no offset column
%          (D-186 of 2026-09-10), stack rebuilt per epoch, coefficient per objective and differentiated through.
%  IV-C 3. M_OBC = M_PS2 with the corrected transition, same criterion, optimiser, selection; THESIS-RESULTS
%          step 5 terms.
%  IV-D 4. Recovery condition as a proposition (own; counterpart of Gyorok 2026 Cond. 4, Thm. 7, Eq. 21);
%          scope: one-step surrogate, aggregate, first order, normalised metric, xbar = 0 slice, offset caveat,
%          Thms. 16, 18 do not transfer.
% CYCLE 11 NOTES (header only)
% - Structure: four subsections (skill Order: problem, component, identification problem, guarantee). New labels
%   sec:obc_problem, eq:obc_sens, eq:obc_tuples, eq:obc_coef, eq:obc_corr, eq:obc_problem, eq:obc_eps. Kept
%   labels sec:obc, sec:obc_negation, sec:obc_map, sec:obc_scope. Removed labels eq:obc_negation,
%   eq:obc_gantry_sens, eq:obc_gyorok, eq:obc_ours, eq:obc_model, eq:obc_delta, eq:obc_bias, eq:obc_overlap.
% - Code confirmed this session: GD/obc_gantry.py (build_reference_set: train_files, k = 2..N-3, q = P^-T y,
%   fourth-order difference, x_a = 0, recorded u_total, normalised with norm.x_mean/std_x; GantryOBCCorrection:
%   subtracts E J(vbar, x_b, u) theta on the routed physical rows, zero on x_a rows; make_basis_builder: jacfwd of
%   transition_from_free over the ten FREE coordinates); FS/obc.py (OBCBasis SVD and pseudoinverse coefficient,
%   OBCLifecycle refresh per epoch, coefficient once per objective with autograd); FS/interconnect.py (loss:
%   maybe_refresh then coefficient; _sync_prediction_state_for_validation: coefficient re-solved before validation,
%   basis rebuilt after restoring the selected checkpoint); FS/lbfgs_polish.py has no OBC logic, so the polish keeps
%   the basis of the restored Adam iterate.
% - Own claims checked by algebra this session: range of the xi-sensitivity equals that of the theta_base
%   sensitivity (d theta_base / d xi triangular, diagonal theta_j and d m_Delta/d xi_Delta nonzero); Proposition 1
%   (residual splits under Pi, no exact-fit assumption needed); Phi^T (I - Pi) = 0 for any rank, so its
%   theta_aug-derivative vanishes; acceleration sensitivity of eq:eom (mass combinations in M only, c_g1, c_g2, c_y,
%   c_bSigma in C only, k_bSigma in K only); fourth-order versus second-order difference gain at 230 Hz, 4 kHz.
% - Corrections to sources: the old draft cited Gyorok 2026 Eq. (39) for parameter uniqueness; the citation log
%   says Eq. (39) does not support it, so Assumption 1 is cited. The old pointwise no-go claim ("no nonzero missing
%   force is orthogonal at a single tuple") was dropped: not checked for every tuple (a tuple at rest loses the
%   damping columns).
% - Resolved original todos: "Source" (citations in the opening); "Subsection title improve" (new titles);
%   "From Section III ... negation" (IV-A cites Kessels and opens from the uniqueness need); "affine arm stricter
%   definition or sensitivity analysis" (neither: the thesis block pins obc_space = tangent, D-200 arm never in the
%   campaign); refs.bib note (gyorok2026obc added, gyorok2025l4dc placeholder replaced, both with % CHECK);
%   "one-step versus trajectory penalty" (one-step: D-192; the trajectory method was ruled out, D-190 Ruled out 1);
%   "orthogonal multisine design" (not run, THESIS-RESULTS; stated as Gyorok Remark 5 in IV-D). Kept and rewritten:
%   Gamma exclusion confirmation by the supervisors.
% - Left out (reader test): rank-loss fallback and abort, float64 solve, hand-written forward sensitivity (D-194),
%   per-objective hoist (D-195), validation monitoring (D-201), Delta* null test (D-203 measurement), overlap
%   measures (not reported by THESIS-RESULTS, todo). Figure: the planned geometric figure does not exist yet.
% - Stale in other files (not edited): 06_results.tex line 40 cites eq:obc_overlap (removed; no overlap display);
%   07_discussion.tex lines 37 and 44 and 99_appendix.tex line 171 write J and Delta* (now Phi and epsilon*; Delta is
%   the scheduling block); 07_discussion.tex line 37 claims pointwise orthogonality is impossible, which this
%   section no longer states; 03_augmentation.tex line 676 promises an estimate that "is unique", which this
%   section delivers only to first order on the tuples (candidate: "unique to first order").
% - Length: about 2.9 pages in a draftfalse build (scratchpad), against 1.55, over 1.5 times after the second
%   deletion pass. Cause: the displays the Math standard requires (sensitivity, tuples, coefficient stack,
%   correction, problem display with the substituted line, discrepancy) and Proposition 1 with its proof, all younger
%   than the budget. Proposed budget: 2.8 pages.
% - Moved out earlier and still out: pointwise no-go and Coulomb inner-product argument (Section VII bullets); rank,
%   condition number and column correlations (Results); dropped rho* = 0.205 and the 4.96 % floor (older data).
\section{Interpretability by Orthogonal Construction}\label{sec:obc}
\ifdraft
\begin{itemize}
\item Opening: from the non-unique split of Section~III-D to a unique estimate of $\theta_{\mathrm{base}}$; names IV-A to IV-D; third contribution of Section~I (style profile; Introduction)
\end{itemize}
\fi

Given the non-unique split of Section~\ref{sec:aug_closed_loop}, our goal is
an estimate of $\theta_{\mathrm{base}}$ that the network cannot take over.
To this end, we state how the network negates parameter changes
(Section~\ref{sec:obc_negation}), construct a correction that removes this
overlap (Section~\ref{sec:obc_map}), identify the corrected model
(Section~\ref{sec:obc_problem}) and delimit its guarantee
(Section~\ref{sec:obc_scope}).
Together they deliver the third contribution of Section~I: state-level
orthogonality for the identifiable physical-parameter tangent space and its
true-parameter recovery condition.
\todo{Missing: contribution wording. The Introduction says ``quantify'' the recovery condition, while this section states the condition and Section~VI reports no value of it. Candidate: ``Construct state-level orthogonality for the identifiable physical-parameter tangent space of a closed-loop, self-scheduled model and state its true-parameter recovery condition'' (THESIS-RESULTS step~5).}

\subsection{Negation of parameter changes}\label{sec:obc_negation}
\ifdraft
\begin{itemize}
\item A unique estimate requires that the network cannot reproduce a change of $\theta_{\mathrm{base}}$, which \eqref{eq:aug_transition} does not ensure (Section~III-D)
\item Parameter directions = sensitivity of the normalised transition in the ten training coordinates $\xi$; raw parameters give rank ten only (Section~II-C; Gy\"or\"ok 2026 Assumption~1)
\item To first order, $f^{\mathrm n}_{\mathrm{aug}}-\Phi v$ cancels the move $\xi\to\xi+v$: Kessels' negation made precise; linear examples of Gy\"or\"ok (Kessels Sec.~5.3.1.3; Gy\"or\"ok 2026 Ex.~2, 2025 Ex.~1; this work)
\item On the gantry a negating write is an inertial, damping or stiffness force, its direction set by $M(Y)^{-1}$ (this work, from \eqref{eq:eom})
\end{itemize}
\fi

A unique estimate of $\theta_{\mathrm{base}}$ requires that the network cannot
reproduce a change of $\theta_{\mathrm{base}}$.
The augmented model \eqref{eq:aug_transition} places no such restriction on
the network.

We measure parameter changes in the training coordinates $\xi$ of
\eqref{eq:mdelta_map}, which parameterise the ten identifiable combinations.
The fourteen raw parameters would give a sensitivity of rank ten only, because
the model does not change along the four lost directions of
Section~\ref{sec:params}.
The parameter directions of the normalised transition are the columns of
\begin{equation}
 \Phi(\xi;\tilde x^{\mathrm n},u^{\mathrm n})
 =\frac{\partial f^{\mathrm n}_{\mathrm{base}}
   \bigl(\theta_{\mathrm{base}}(\xi),\tilde x^{\mathrm n},u^{\mathrm n}\bigr)}
   {\partial\xi},
 \label{eq:obc_sens}
\end{equation}
where $\Phi:\R^{10}\times\R^6\times\R^3\to\R^{6\times10}$.
To first order in a displacement $v\in\R^{10}$, replacing $\xi$ by $\xi+v$
and $f^{\mathrm n}_{\mathrm{aug}}$ by
$f^{\mathrm n}_{\mathrm{aug}}-\Phi(\xi;\tilde x^{\mathrm n},u^{\mathrm n})v$
leaves the transition of \eqref{eq:aug_transition} unchanged.

This first-order cancellation makes the negation of Kessels, learned terms
that partially negate the first-principles terms, precise for the parameter
directions~\cite[Sec.~5.3.1.3]{kessels2025ai}.
Gy\"or\"ok et al.\ show the same non-uniqueness for linear baselines
augmented with a linear layer~\cite[Ex.~2]{gyorok2026obc},
\cite[Ex.~1]{gyorok2025l4dc}.

On the gantry, a negating write acts as a force on the inertia.
The columns of $\Phi$ propagate through the RK4 step the acceleration
sensitivity $-M(Y)^{-1}\bigl(\partial_jM(Y)\,\ddot q+\partial_jC\,\dot q
+\partial_jK\,q\bigr)$ of \eqref{eq:eom}, with $\partial_j$ the derivative
with respect to $\xi_j$.
The inertial combinations $m_h$, $m_\Sigma$, $m_\Delta$, $J_{\mathrm{eff}}$
and $d$ therefore act through $\ddot q$.
The damping combinations act through $\dot q$.
The stiffness $k_{b,\Sigma}$ acts through $q$.
The factor $M(Y)^{-1}$ makes each direction depend on the payload position.

\subsection{Orthogonal-by-construction correction}\label{sec:obc_map}
\ifdraft
\begin{itemize}
\item The network must stay out of $\ran\Phi$ on a transition that is rational in $\theta_{\mathrm{base}}$, with only positions recorded; Gy\"or\"ok 2026 names this state-space setting future work (Sec.~6)
\item Reference tuples from every training sample: $q=P^{-\top}y$, fourth-order velocity difference (band gain, damping multiplies velocity), $\bar x=0$, recorded input (Gy\"or\"ok 2026 Sec.~3.2 auxiliary set; code; implementation plan; reasons (?))
\item Adapted: least-squares coefficient on the stacked sensitivity at $\xi^\circ$, no weight, orthogonal stack (Gy\"or\"ok 2026 Eqs.~(8), (10), Lemma~3; 2025 Sec.~4); offset not protected (D-186; supervisors (?))
\item Correction at the rollout point with the fixed coefficient (Gy\"or\"ok 2026 Eq.~(13)); physical rows only, $g_{\mathrm{aug}}$ uncorrected (D-192)
\item Consequence: to first order the stacked one-step map determines $\xi$ (this work)
\item Stack per epoch at the current $\xi$, coefficient per evaluation of $V$, differentiated through; gradient to $\xi$ only through $f^{\mathrm n}_{\mathrm{base}}$ (D-190 items~5, 6; D-192; Gy\"or\"ok 2025 Sec.~4, Remark~2; measured reason (?))
\end{itemize}
\fi

The network must stay out of $\ran\Phi$, on a transition that is rational in
$\theta_{\mathrm{base}}$ and whose states are only partly recorded.
Gy\"or\"ok et al.\ name a state-space baseline without full-state
measurement as future work for their method~\cite[Sec.~6]{gyorok2026obc}.
They allow the construction to be evaluated on any auxiliary set of
regressor evaluations~\cite[Sec.~3.2]{gyorok2026obc}.
The records hold positions and actuator forces, but no velocities and no
added states.

We therefore build reference tuples from every sample of the training records
at which a fourth-order central difference is defined, because the damping
directions of $\Phi$ multiply velocity.
Its gain error grows with the fourth power of frequency, against the square
for a second-order difference:
\begin{equation}
\begin{aligned}
 \tilde x^{\mathrm n}_i&=T_x\bigl(\col(q_i,\dot q_i)-\mu_x\bigr),
 \qquad q_i=P^{-\top}y_i,\\
 \hat x^{\mathrm n}_i&=\col(\tilde x^{\mathrm n}_i,0),
 \qquad u^{\mathrm n}_i=T_u(u_{\mathrm{act},i}-\mu_u),
 \qquad i\in\mathcal R,
\end{aligned}
\label{eq:obc_tuples}
\end{equation}
where $y_i$ and $u_{\mathrm{act},i}$ are the recorded positions and actuator
forces of sample $i$, $\dot q_i$ is the central difference of $q$ over the
samples $i-2$ to $i+2$, and $\mathcal R$ is the set of these samples, with
$N_{\mathcal R}$ elements.
The recorded forces already contain the feedback, so the controller does not
enter the tuples.
\todo{Missing: reason for tuples from the records rather than from a forward simulation of the baseline, as Gy\"or\"ok et al.\ do without state measurements~\cite[Sec.~3]{gyorok2025l4dc}. Candidate: an open-loop simulation leaves position offsets on the free axes (Section~\ref{sec:aug_loop}), so its states would leave the records (own reading; D-192 gives no reason).}
\todo{Missing: reason for $\bar x=0$ in the tuples. D-192 item~2 fixes it without a reason. Candidate: the records fix no value of $\bar x$, which is defined only up to a change of coordinates (Section~\ref{sec:aug_structure}; own reading).}
\todo{Missing: the velocities of \eqref{eq:aug_norm} are first differences, those of the tuples fourth-order differences; no decision records why the two differ. Candidate: the tuples feed the parameter directions, the normalisation only scales (OBC implementation plan in the source-map documentation; own reading).}

For a baseline linear in its parameters, Gy\"or\"ok et al.\ fit the stacked
network output to the stacked regressor in least squares and subtract the
fit~\cite[Eqs.~(8), (10)]{gyorok2026obc}.
Their earlier penalty traded orthogonality against fit through a
weight~\cite[Eq.~(13), Remark~1]{gyorok2025l4dc}.
The construction needs no such weight~\cite[Sec.~3.2]{gyorok2026obc}.
For a baseline nonlinear in its parameters, they take the regressor as the
parameter sensitivity at an expansion point and extend it with the offset of
that expansion~\cite[Sec.~4, Eqs.~(15) to (18)]{gyorok2025l4dc}.
We protect the parameter directions only.
Protecting the offset would charge a direction along which the estimate of
$\xi$ does not move.
It would also tie the protected range to the origin of the state
coordinates, which the centring of \eqref{eq:aug_norm} sets.
\todo{Missing: supervisors' confirmation of the tangent range. D-190 recorded their condition as non-overlap with the baseline output, which is the affine range; D-192 and D-186 (2026-09-10) replace it by the parameter directions. Candidate: keep the tangent range with the two reasons above.}

On the tuples, with the sensitivity \eqref{eq:obc_sens} at an expansion point
$\xi^\circ$ as regressor, the construction becomes
\begin{equation}
\begin{aligned}
 \boldsymbol\Phi(\xi^\circ)&=\col\bigl(\Phi(\xi^\circ;\tilde x^{\mathrm n}_i,
   u^{\mathrm n}_i)\bigr)_{i\in\mathcal R},\\
 \boldsymbol f^{\mathrm n}_{\mathrm{aug}}(\theta_{\mathrm{aug}})
 &=\col\bigl(f^{\mathrm n}_{\mathrm{aug}}(\theta_{\mathrm{aug}},
   \hat x^{\mathrm n}_i,u^{\mathrm n}_i)\bigr)_{i\in\mathcal R},\\
 \theta_{\mathrm{aux}}(\theta_{\mathrm{aug}})
 &=\boldsymbol\Phi(\xi^\circ)^{+}\,
   \boldsymbol f^{\mathrm n}_{\mathrm{aug}}(\theta_{\mathrm{aug}}),\\
 0&=\boldsymbol\Phi(\xi^\circ)^\top\bigl(
   \boldsymbol f^{\mathrm n}_{\mathrm{aug}}(\theta_{\mathrm{aug}})
   -\boldsymbol\Phi(\xi^\circ)\,\theta_{\mathrm{aux}}(\theta_{\mathrm{aug}})\bigr),
\end{aligned}
\label{eq:obc_coef}
\end{equation}
where $\xi^\circ\in\R^{10}$,
$\boldsymbol\Phi(\xi^\circ)\in\R^{6N_{\mathcal R}\times10}$ and
$\boldsymbol f^{\mathrm n}_{\mathrm{aug}}(\theta_{\mathrm{aug}})
\in\R^{6N_{\mathcal R}}$ are the stacks,
$\theta_{\mathrm{aux}}:\R^{n_{\mathrm{aug}}}\to\R^{10}$ is the coefficient,
and $(\cdot)^{+}$ is the Moore-Penrose pseudoinverse.
The last line is Lemma~3 of~\cite{gyorok2026obc} and holds for every
$\theta_{\mathrm{aug}}$.
\todo{Missing: the numerical rank of $\boldsymbol\Phi(\xi^\circ)$ on the thesis tuples; THESIS-RESULTS does not report it, while the full-rank assumption of Proposition~1 needs it. Candidate: one sentence in Section~\ref{sec:results} with the rank at the first and the last rebuild.}

At a rollout point, the correction uses this coefficient with the current
state and input, as in the prediction of Gy\"or\"ok et
al.~\cite[Eq.~(13)]{gyorok2026obc}.
It acts on the six physical rows only, since the baseline transition has no
added-state rows:
\begin{equation}
\begin{aligned}
 f^{\perp}_{\mathrm{aug}}&(\theta_{\mathrm{aug}},\hat x^{\mathrm n}_k,u^{\mathrm n}_k)
 =f^{\mathrm n}_{\mathrm{aug}}(\theta_{\mathrm{aug}},\hat x^{\mathrm n}_k,u^{\mathrm n}_k)\\
 &-\Phi(\xi^\circ;\tilde x^{\mathrm n}_k,u^{\mathrm n}_k)\,
  \theta_{\mathrm{aux}}(\theta_{\mathrm{aug}}),
\end{aligned}
 \label{eq:obc_corr}
\end{equation}
where $f^{\perp}_{\mathrm{aug}}:\R^{n_{\mathrm{aug}}}\times\R^{6+n_{\bar x}}
\times\R^3\to\R^6$ is the corrected learned write.
The added-state update $g_{\mathrm{aug}}$ stays uncorrected.

The correction thereby removes the negation of
Section~\ref{sec:obc_negation} on the tuples.
At the tuples, $f^{\perp}_{\mathrm{aug}}$ stacks to the orthogonal vector of
the last line of \eqref{eq:obc_coef}.
A change $v$ of $\xi$ near $\xi^\circ$ moves the stacked baseline transition
by $\boldsymbol\Phi(\xi^\circ)v$ to first order.
The corrected write stays orthogonal to $\ran\boldsymbol\Phi(\xi^\circ)$ for
every $\theta_{\mathrm{aug}}$, so it cannot reproduce that move.
To first order, the stacked one-step map therefore determines $\xi$.
Since $\partial\theta_{\mathrm{base}}/\partial\xi$ is nonsingular in
\eqref{eq:mdelta_map}, the same holds for $\theta_{\mathrm{base}}$.

The stack and the coefficient are rebuilt at different rates.
At every epoch boundary, $\xi^\circ$ is set to the current $\xi$.
$\boldsymbol\Phi(\xi^\circ)$ is then rebuilt and held fixed for the epoch.
Gy\"or\"ok et al.\ update the expansion point at every objective evaluation or
fix it at the nominal parameters~\cite[Sec.~4]{gyorok2025l4dc}.
Updating rebuilds the full stack at every evaluation.
A fixed point loses accuracy as the estimate moves away from it.
The per-epoch rebuild adapts their recomputation of the projection at the
start of each epoch, which refreshes state estimates rather than the
expansion point~\cite[Remark~2]{gyorok2025l4dc}.
The coefficient $\theta_{\mathrm{aux}}$ is recomputed from the current
network at every evaluation of the criterion.
Gradients reach $\theta_{\mathrm{aug}}$ through $f^{\mathrm n}_{\mathrm{aug}}$
and through $\theta_{\mathrm{aux}}$, and reach $\xi$ only through
$f^{\mathrm n}_{\mathrm{base}}$.
The last line of \eqref{eq:obc_coef} holds for every $\theta_{\mathrm{aug}}$,
so its derivative with respect to $\theta_{\mathrm{aug}}$ vanishes too.
An update of the network therefore never moves the corrected stack into
$\ran\boldsymbol\Phi(\xi^\circ)$.
\todo{Missing: the reason for the epoch rate is a measurement: the sensitivity range turns slowly with $\xi$ (D-190 item~6), which no section reports. Candidate: report the principal angle between the first and the last stack in Section~\ref{sec:results}.}

\subsection{Identification of the corrected model}\label{sec:obc_problem}
\ifdraft
\begin{itemize}
\item The corrected model is identified with the criterion, optimiser and selection of $\mathcal M_{\mathrm{PS2}}$; trained vector unchanged, $\theta_{\mathrm{aux}}$ no parameter (THESIS-RESULTS Sec.~2; code)
\item Problem display: cost of \eqref{eq:aug_loss}, changed transition line with the correction substituted, the rest cited (Hoekstra 2026 Eq.~(22); README Math standard item~3)
\item Adam, L-BFGS polish and PS2 as in Section~III; the polish keeps the stack of the restored Adam iterate (code)
\item Selection as in Section~III-D; coefficient recomputed before validation; stack rebuilt at the selected parameters (D-198; code)
\item $\mathcal M_{\mathrm{OBC}}$ named; compared with $\mathcal M_{\mathrm{PS2}}$ from both starts in THESIS-RESULTS step~5 terms (notation (?))
\end{itemize}
\fi

To compare the corrected model with $\mathcal M_{\mathrm{PS2}}$, we identify
it with the same criterion, optimiser and selection.
The trained parameters remain
$\vartheta=(\xi,\theta_{\mathrm{aug}},\theta_{\mathrm{enc}})$ of
Section~\ref{sec:aug_closed_loop}.
The coefficient $\theta_{\mathrm{aux}}$ is not trained, as it is a function of
$\theta_{\mathrm{aug}}$.
With \eqref{eq:obc_corr} substituted, the identification problem is
\begin{equation}
\begin{aligned}
 &\min_{\vartheta}\;V(\vartheta)\quad\text{s.t.}\\
 &\tilde x^{\mathrm n}_{k+1}
 =f^{\mathrm n}_{\mathrm{base}}\bigl(\theta_{\mathrm{base}}(\xi),\tilde x^{\mathrm n}_k,\hat u^{\mathrm n}_k\bigr)
  +f^{\mathrm n}_{\mathrm{aug}}(\theta_{\mathrm{aug}},\hat x^{\mathrm n}_k,\hat u^{\mathrm n}_k)\\
 &\qquad\quad-\Phi(\xi^\circ;\tilde x^{\mathrm n}_k,\hat u^{\mathrm n}_k)\,
  \theta_{\mathrm{aux}}(\theta_{\mathrm{aug}}),
 \qquad\theta_{\mathrm{aux}}\text{ from }\eqref{eq:obc_coef},\\
 &\text{the remaining constraints of }\eqref{eq:aug_loss},
\end{aligned}
\label{eq:obc_problem}
\end{equation}
where $V$ is the criterion of \eqref{eq:aug_loss}.

Adam, the L-BFGS polish and the PS2 stages of Section~\ref{sec:aug_ps2} run as
for $\mathcal M_{\mathrm{PS2}}$, with this transition.

The parameters are selected as in Section~\ref{sec:aug_closed_loop}.
Before every validation run, $\theta_{\mathrm{aux}}$ is recomputed from the
network as the last update left it.
After the selected Adam iterate is restored, $\boldsymbol\Phi$ is rebuilt at
its $\xi$ and $\theta_{\mathrm{aux}}$ is recomputed.
The polish keeps this stack.

We call the model \eqref{eq:aug_transition} with
$f^{\mathrm n}_{\mathrm{aug}}$ replaced by $f^{\perp}_{\mathrm{aug}}$,
identified by \eqref{eq:obc_problem} from $\theta_{\mathrm{base},0}$ with PS2,
and selected as above, $\mathcal M_{\mathrm{OBC}}$.
Its inputs and output are those of $\mathcal M_{\mathrm U}$.
Section~\ref{sec:results} compares $\mathcal M_{\mathrm{OBC}}$ with
$\mathcal M_{\mathrm{PS2}}$ from the two starts of Section~\ref{sec:setup}, to
show what the correction does to the ten combinations and at what cost in
accuracy.
\todo{Missing: model notation. Candidate: $\mathcal M_{\mathrm{OBC}}$, calligraphic as $\mathcal M_{\mathrm{LPV}}$ and $\mathcal M_{\mathrm{PS2}}$, since $M$ is the inertia matrix (README open conflict on model names).}

\subsection{Scope and condition for true parameter recovery}\label{sec:obc_scope}
\ifdraft
\begin{itemize}
\item A unique estimate is not the true one: the part of the omitted dynamics inside $\ran\boldsymbol\Phi$ is still attributed to $\xi$ (Gy\"or\"ok 2026 Sec.~4.1)
\item True one-step discrepancy on the tuples with a successor, same difference stencil (D-202, D-203)
\item Proposition: under local affinity, full rank and a one-step least-squares fit, $\hat\xi=\xi^*+\boldsymbol\Phi^{+}\boldsymbol\varepsilon^*$, recovery iff $\boldsymbol\Phi^\top\boldsymbol\varepsilon^*=0$ (this work; counterpart of Gy\"or\"ok 2026 Cond.~4, Thm.~7, Eq.~(21))
\item Designed excitation can meet the condition (Gy\"or\"ok 2026 Remark~5); whether the records meet it is not reported (?)
\item Scope: one-step surrogate of $V$; aggregate, first order; normalised metric; $\bar x=0$ slice and uncorrected $g_{\mathrm{aug}}$, active under PS2; offset caveat; noise in $\boldsymbol\varepsilon^*$; Thms.~16, 18 do not transfer (D-186; D-192; THESIS-RESULTS step~5)
\end{itemize}
\fi

Even a unique estimate is not the true one, because the part of the omitted
dynamics inside $\ran\boldsymbol\Phi$ is still attributed to $\xi$.
Let $\xi^*$ give the true combinations,
$\theta_{\mathrm{base}}(\xi^*)=\theta^*_{\mathrm{base}}$.
On the tuples whose successor is also a tuple, the true one-step discrepancy
is
\begin{equation}
 \varepsilon^*_i=\tilde x^{\mathrm n}_{i+1}
 -f^{\mathrm n}_{\mathrm{base}}\bigl(\theta_{\mathrm{base}}(\xi^*),
   \tilde x^{\mathrm n}_i,u^{\mathrm n}_i\bigr),
 \qquad i\in\mathcal R_+,
 \label{eq:obc_eps}
\end{equation}
where $\mathcal R_+\subset\mathcal R$ holds these tuples and
$\varepsilon^*_i\in\R^6$, stacked as $\boldsymbol\varepsilon^*$.
The successor $\tilde x^{\mathrm n}_{i+1}$ is formed with the same difference
stencil, so the velocity rows of $\varepsilon^*_i$ contain no difference of
two stencils.

\textit{Proposition 1:} Let $\boldsymbol\Phi$, $\boldsymbol\varepsilon^*$,
the stacked baseline transitions $\boldsymbol f^{\mathrm n}_{\mathrm{base}}(\xi)$
and the stacked successors $\tilde{\boldsymbol x}^{\mathrm n}_{+}$ run over
$\mathcal R_+$, and let $\Pi=\boldsymbol\Phi\boldsymbol\Phi^{+}$.
Suppose
(i) on a set $\Xi\subset\R^{10}$ that contains $\xi^*$ and
$\xi^*+\boldsymbol\Phi^{+}\boldsymbol\varepsilon^*$,
$\boldsymbol f^{\mathrm n}_{\mathrm{base}}$ is affine in $\xi$ with slope
$\boldsymbol\Phi$, which the gantry transition meets only locally;
(ii) $\boldsymbol\Phi$ has full column rank;
(iii) $(\hat\xi,\hat\theta_{\mathrm{aug}})$ minimises
$\norm{\tilde{\boldsymbol x}^{\mathrm n}_{+}
-\boldsymbol f^{\mathrm n}_{\mathrm{base}}(\xi)
-(I-\Pi)\boldsymbol f^{\mathrm n}_{\mathrm{aug}}(\theta_{\mathrm{aug}})}_2^2$
over $\Xi\times\R^{n_{\mathrm{aug}}}$.
Then $\hat\xi=\xi^*+\boldsymbol\Phi^{+}\boldsymbol\varepsilon^*$, and
$\hat\xi=\xi^*$ if and only if
$\boldsymbol\Phi^\top\boldsymbol\varepsilon^*=0$.

\begin{IEEEproof}
By (i) and \eqref{eq:obc_eps}, the residual of (iii) is
$\boldsymbol\Phi(\xi^*-\xi)+\boldsymbol\varepsilon^*
-(I-\Pi)\boldsymbol f^{\mathrm n}_{\mathrm{aug}}$.
Its components in $\ran\boldsymbol\Phi$ and in the orthogonal complement are
$\boldsymbol\Phi(\xi^*-\xi)+\Pi\boldsymbol\varepsilon^*$ and
$(I-\Pi)(\boldsymbol\varepsilon^*-\boldsymbol f^{\mathrm n}_{\mathrm{aug}})$.
The squared norm is the sum of their squared norms, and only the first depends
on $\xi$.
By (ii), the first vanishes only at
$\xi=\xi^*+\boldsymbol\Phi^{+}\boldsymbol\varepsilon^*$, which lies in $\Xi$
by (i).
Every minimiser in (iii) therefore has this $\hat\xi$.
By (ii), $\boldsymbol\Phi^{+}=(\boldsymbol\Phi^\top\boldsymbol\Phi)^{-1}
\boldsymbol\Phi^\top$, so $\boldsymbol\Phi^{+}\boldsymbol\varepsilon^*=0$ if
and only if $\boldsymbol\Phi^\top\boldsymbol\varepsilon^*=0$.
\end{IEEEproof}

Proposition~1 is the state-level counterpart of Condition~4 and the error
(21) of Gy\"or\"ok et al.~\cite[Sec.~4.1, Thm.~7]{gyorok2026obc}.
It needs no exact fit of the data, their Assumption~6, because the two
components of the residual separate.
Gy\"or\"ok et al.\ note that an input designed from the known structure of the
omitted term can meet the condition~\cite[Remark~5]{gyorok2026obc}.
\todo{Missing: whether the training records meet $\boldsymbol\Phi^\top\boldsymbol\varepsilon^*\approx0$. THESIS-RESULTS step~5 states that the omitted physics of the truth is not orthogonal, but its only evidence is the overlap measured on older data (D-190). Candidate: report $\norm{\Pi\boldsymbol\varepsilon^*}/\norm{\boldsymbol\varepsilon^*}$ on the thesis tuples in Section~\ref{sec:results}, labelled as using the true combinations.}

The construction and Proposition~1 hold within the following scope.
Proposition~1 concerns a one-step least-squares fit, whereas
$\mathcal M_{\mathrm{OBC}}$ minimises the closed-loop multistep criterion
$V$.
It is therefore a surrogate for the trained estimate, not its bias.
The orthogonality holds for the stack over the tuples, not per tuple.
Its consequence for $\xi$ holds to first order at $\xi^\circ$.
Their inner product is that of the normalised coordinates, so $T_x$ weights
the state rows.
The tuples lie at $\bar x=0$, so $\theta_{\mathrm{aux}}$ sees no direction in
which $f^{\mathrm n}_{\mathrm{aug}}$ acts through nonzero added states.
With PS2 these states are active.
The uncorrected $g_{\mathrm{aug}}$ reaches the physical rows one step later
through $f^{\mathrm n}_{\mathrm{aug}}$.
The construction therefore implies no orthogonality of simulated trajectories
or outputs.
A learned write that cancels the nominal transition
$f^{\mathrm n}_{\mathrm{base}}$ itself is constrained only through its
component in $\ran\boldsymbol\Phi(\xi^\circ)$.
If the records are noisy, $\boldsymbol\varepsilon^*$ also contains the noise
and the error of the difference stencil.
The consistency and zero-covariance results of Gy\"or\"ok et
al.~\cite[Thms.~16, 18]{gyorok2026obc} assume a linear-in-parameters
input-output baseline and do not transfer.
